\section{A direct construction with twelve positive unknowns}
\label{sec:radical-twelve}

The nine-variable compiler and the smaller arithmetic circuits optimize
different measures.  Here is a further direct reduction
of the native/Pell family: twelve strictly positive unknowns, three
equations of maximum degree $1600$, or one equation of degree $1609$.
The emitted single-equation evaluator has $151$ binary operations
($97$ multiplications and $54$ additions or subtractions); the system
uses $147$.  Thus it improves the thirteen-variable direct construction
in variable count at the cost of degree and arithmetic size.
All four coefficient ports below are bound together to the fixed T4 tuple
specified earlier.  No ports from different compiler recipes are mixed.

This section uses two established interface facts from the source
packages.  Package~09 proves its weak-sign bootstrap and weighted-root
inequality for the native four-port interface.  Package~14 proves the
signed-core soundness argument, including recovery of the main-root sign.
In that soundness argument the old strong multiplier need only be nonzero
and the scale $w$ positive.  Oddness of that multiplier and squareness of
$w$ are sufficient conditions used in its particular unrestricted-integer
parameterization, and are not additional hypotheses on the positive
construction below.  These arithmetic interface arguments transfer with
the whole T4 tuple.

\subsection{The polynomial}

The unknowns are
\[
 J,F,\alpha,z,h,i,T,s,w,\eta,\zeta,t\in\mathbb Z_{>0}.
\]
All other symbols in this subsection abbreviate polynomials.  Write
$\beta=B-1$ and let $K_0,M_C,M_F$ be the other three fixed ports,
with the shifted $M_F$ convention.  Define
\begin{align*}
 q&=\beta J+1,& C_w&=q-F-\alpha,& X&=wq,&Y&=sq^3,\\
 E&=XY,&k&=\eta+\zeta,&c&=kY+\eta,&a&=E+Y,\\
 H&=4a+3,&\Delta&=a^2+H=(a+2)^2-1.
\end{align*}
\begin{align*}
 R&=((q-1)(q-F)+\alpha)(q^2-1)+(M_C+qM_F)J,\\
 N_k&=k-hE-R,&
 N_t&=(K_0+w)C_w+q-F-z(q-1),\\
 L_0&=EkY,& A&=1+L_0(L_0+k).
\end{align*}
The new strong multiplier is the old multiplier restricted to $Hi$:
\begin{align*}
 I&=Hic^2,& D&=1+\Delta I^2,&Q&=\Delta^2I^2,&d&=Q-1,\\
 v&=cT,&z_0&=c+RD,& e&=Q(Dv^2+z_0^2)-1,&b&=2Qvz_0,\\
 M&=2(D+v^2),&G&=1+\Delta c^2.
\end{align*}
For the separating weight put
\begin{align*}
 T_*&=L_0+k+1,&F_*&=(a+2)I+1,\\
 Y_*&=2(vF_*+z_0),&U_*&=MT_*+Y_*,\\
 \ell&=1+HU_*,&S&=Ht+\ell(ac+X),&W&=HM.
\end{align*}

The following fixed computation defines a polynomial without introducing
an unknown.  Work formally modulo $r^2-G$ and form
\begin{align*}
 Z&=S-\ell r,\\
 U&=d(Z^2+W^2A)-H^2e,& V&=2dWZ,\\
 F_0&=U^2+AV^2-H^4b^2D,&F_1&=2UV,\\
 F_0^2-AF_1^2&=E_0+E_1r\pmod{r^2-G},\\
 \mathcal R_{12}&=E_0^2-GE_1^2.
\end{align*}
Every reduction replaces $r^2$ by $G$, so $E_0,E_1$ are ordinary
integer polynomials in the twelve supplied unknowns.  This is a finite
arithmetic recipe, not an existential use of $r$.  The supplied JSON
circuit performs it with coefficient pairs; in particular it evaluates
$H^4b^2D$ as $4H^4(QDv^2)(Qz_0^2)$.

Our system and single equation are respectively
\begin{equation}
 N_k=1,\qquad N_t=1,\qquad\mathcal R_{12}=0,
 \label{eq:twelve-system}
\end{equation}
and
\begin{equation}
 \boxed{\mathcal P_{12}=N_kN_t(1+\mathcal R_{12})-1=0.}
 \label{eq:twelve-product}
\end{equation}

\begin{theorem}
With the stated fixed coefficient binding, each of
\eqref{eq:twelve-system} and \eqref{eq:twelve-product} has a positive
integer solution exactly when K215 halts on blank tape.  Consequently
ZFC proves that its absence of positive integer solutions is equivalent
to $\operatorname{Con}(\mathrm{ZFC})$.
\end{theorem}

\subsection{Soundness: separating and recovering four roots}

We prove soundness under the more general assumptions
$N_k=N_t=\varepsilon\in\{1,-1\}$ and $\mathcal R_{12}=0$;
this also covers the product finalizer.  The inherited weak-sign
bootstrap gives
\[
 q\ge16,\quad 0<R<q^4-2,\quad E\ge q^4>R+2,\quad k>R,
\]
and, with $M_0=D+v^2$, the inequality
\[
 dM_0^2A>Q(Dv^2+z_0^2)=e+1.
\]
These steps do not use the main root or its discarded coordinate
$\gamma$.  Since $M=2M_0$, they imply
\begin{equation}
 dM^2A>4(e+1),\qquad b\sqrt D\le e+1.
 \label{eq:twelve-domination}
\end{equation}
The second inequality is $2v z_0\sqrt D\le Dv^2+z_0^2$.
All quantities used as positive bounds are positive, with $d\ge1$ and
$D,G,A>1$.

For proof purposes, $\mathcal R_{12}$ is the product over the eight
choices $r=\pm\sqrt G$, $f=\pm\sqrt D$, and
$\tau=\pm\sqrt A$ of
\begin{equation}
 d(S-\ell r-HM\tau)^2-H^2(e-bf).
 \label{eq:twelve-branch-factor}
\end{equation}
Taking the quadratic norms first in $f$, then in $\tau$, then in $r$
gives exactly the displayed polynomial computation.  Its zero supplies
a real algebraic tuple satisfying
\begin{equation}
 r^2=G,\quad f^2=D,\quad\tau^2=A,\quad
 dy^2=e-bf,\quad S=\ell r+H(M\tau+y).
 \label{eq:twelve-radicals}
\end{equation}
Indeed the first three radicals are real, and the last equality defines
a real $y$.

There is at most one such tuple for the given rational $S$.  First,
$\sqrt A<T_*$ and $\sqrt D<F_*$.  Moreover
\[
 y^2=\frac{Q(vf-z_0)^2-1}{Q-1}
       <2(v\sqrt D+z_0)^2,
 \qquad |y|<Y_*.
\]
It follows that $|M\tau+y|<U_*$.  Two different signs of $r$ would
therefore change $\ell r$ by $2\ell\sqrt G>2HU_*$, more than the
remaining contribution can change.  Thus $r$ is unique.  Next,
\eqref{eq:twelve-domination} gives
\[
 M^2 A>\frac{e+b\sqrt D}{d}\ge y^2.
\]
Two different signs of $\tau$ cannot give the same $M\tau+y$,
because their difference $2M\sqrt A$ exceeds the possible difference
of the $y$ terms.  Therefore $\tau$, then $y$, then $f$ are unique;
the last conclusion uses $b\ne0$.

Every algebraic embedding of the field generated by this tuple fixes
$S$ and the integer coefficients, and produces another tuple in
\eqref{eq:twelve-radicals}.  Its $r,f,\tau$ are real; the last equality
forces its $y$ to be real as well.  Uniqueness fixes every generator
under every embedding.  Characteristic-zero separability, or the fixed
field theorem in a normal closure, now implies that all four roots are
rational.

The rational numbers $r,f,\tau$ have integer squares and hence are
integers.  The last equality in \eqref{eq:twelve-radicals} makes $Hy$
an integer.  If $y=u/v_0$ is in lowest terms, then $v_0\mid H$ and
$v_0^2\mid d$.  But
\[
 Q=\Delta^2H^2i^2c^4,\qquad \gcd(H,Q-1)=1.
\]
Thus $y$ is an integer.  Set
\[
 \gamma=(r-ac-X)/H.
\]
The same equality yields $t=\ell\gamma+M\tau+y$, so
$\ell\gamma$ is an integer.  Since $\gcd(\ell,H)=1$, $\gamma$ is
an integer too.  We have restored the exact equations
\begin{align*}
 (ac+X+\gamma H)^2-\Delta c^2&=1,\\
 f^2-\Delta(Hi)^2c^4&=1,\\
 \tau^2&=1+L_0(L_0+k),\\
 QV_a^2-(Q-1)y^2&=1,
 \qquad V_a=c(Tf-1)-Rf^2.
\end{align*}
Here $e-bf=Q(vf-z_0)^2-1$ because $f^2=D$.
Also $y\ne0$: otherwise $QV_a^2=1$ would contradict $Q>1$.
Taking absolute values restores positive $\tau,y$.

The signed-core soundness argument cited above applies to the possibly
signed $f,\gamma$ (equivalently signed $T$ after changing the sign of
$f$).  Its normalized rank step uses only the nonzero multiplier $Hi$.
It recovers the main index $p=R$, the positive main-root branch,
$\varepsilon=1$, and the native masks.  Thus it decodes a blank
accepting K215 tableau.  This use of signed-core soundness is essential:
the older theorem restricted to positive $\gamma$ alone would not
justify this step.

\subsection{Completeness with the new divisibility restriction}

Take the canonical positive witnesses supplied by a blank accepting
tableau.  They have $R\equiv3\pmod4$,
$c=\psi_{a+2}(R)$ and main root $\chi_{a+2}(R)$, where
$\chi_A(n)+\psi_A(n)\sqrt{A^2-1}=(A+\sqrt{A^2-1})^n$.
Keep the basic coordinates, main $\gamma$, and $\tau$.  We must
rechoose the strong and auxiliary witnesses; an arbitrary old witness
need not have its multiplier divisible by $H$.

Let $N=Hc^2$.  The matrix
\[
 \begin{pmatrix}
 \chi_{a+2}(R)&\Delta\psi_{a+2}(R)\\
 \psi_{a+2}(R)&\chi_{a+2}(R)
 \end{pmatrix}
\]
has determinant one.  Modulo $N$ it belongs to the finite group
$\mathrm{SL}_2(\mathbb Z/N\mathbb Z)$, whose cardinality is at most
$N^4$.  Its order divides $(N^4)!$.  Consequently, for
$m=R(N^4)!$, the Pell multiplication law gives
$Hc^2\mid\psi_{a+2}(m)$.  Put
\[
 i=\frac{\psi_{a+2}(m)}{Hc^2},\quad
 f=\chi_{a+2}(m),\quad \rho=\Delta\psi_{a+2}(m),\quad
 V_a=\frac{\chi_\rho(R)}\rho,\quad y=\psi_\rho(R).
\]
These are positive integers; oddness of $R$ gives integrality of $V_a$.
The odd Chebyshev quotient identities, using $R\equiv3\pmod4$,
give
\[
 V_a\equiv-c\pmod f,\qquad V_a\equiv-R\pmod c,
 \qquad f^2\equiv1\pmod c.
\]
For clarity, their inputs are
$\rho^2\equiv-\Delta=1-(a+2)^2\pmod f$ and
$\rho^2\equiv0\pmod c$; they do not require the old auxiliary index.
Thus
\[
 o=(V_a+c)/f,\qquad T=(o+Rf)/c
\]
are positive integers.  The second divisibility follows by multiplying
$of\equiv-R\pmod c$ by $f$ and using $f^2\equiv1\pmod c$.
They restore exactly $V_a=c(Tf-1)-Rf^2$.
Finally choose $t=\ell\gamma+M\tau+y>0$.  Then
\eqref{eq:twelve-radicals} holds, proving completeness.

For the single equation, a zero of \eqref{eq:twelve-product} makes
all three factors integer units.  Since $\Delta\equiv0$ or $3\pmod4$,
we have $G\equiv0$ or $1\pmod4$.  Therefore
$\mathcal R_{12}=E_0^2-GE_1^2$ cannot be $2\pmod4$:
the factor $1+\mathcal R_{12}$ cannot be $-1$.  It must be $1$,
and $N_k,N_t$ have equal signs.  The soundness proof deliberately
allowed those signs, and so proves the product claim as well.
Composing with the K215/refutation equivalence proves the theorem.

\subsection{Degree, reproducibility, and further questions}

The degrees relevant to the calculation are
\[
\begin{array}{c|rrrrrrrrrr}
 \text{quantity}&c&a&\Delta&L_0&A&I&D&Q&z_0&e\\ \hline
 \deg&5&6&12&11&22&17&46&58&50&158
\end{array}
\]
and
\[
 \deg b=114,\quad\deg M=46,\quad\deg U_*=57,\quad
 \deg\ell=63,\quad\deg S=74,\quad\deg W=52.
\]
For an exact certificate, replace every unknown by $\lambda$ times
an independent indeterminate and work with formal Laurent series at
infinity.  The two roots of $G$ have expansions
$r_+=ac+2c+\text{terms of degree below }5$ and $r_-=-r_+$.
Thus $Z_+=S-\ell r_+$ has degree $68$, with leading form
$-2\ell_{\rm top}c_{\rm top}$, and $Z_-$ has degree $74$, with
leading form $2\ell_{\rm top}a_{\rm top}c_{\rm top}$.
The term $W\tau$ has degree $63$; the terms $H^2e$ and $H^2bf$
have degrees $170$ and $149$.  Hence the leading term of each factor
\eqref{eq:twelve-branch-factor} is $dZ^2$, of degree $194$ or $206$.
There are four of each.  This proves
\[
 \deg\mathcal R_{12}=4(194)+4(206)=1600,
 \qquad\deg\mathcal P_{12}=1600+7+2=1609.
\]
The leading form of $\mathcal R_{12}$ is the nonzero polynomial
$2^{16}d_{\rm top}^{8}\ell_{\rm top}^{16}
a_{\rm top}^{8}c_{\rm top}^{16}$, so the claim is uniform for valid
fixed coefficients.

The generator emits every gate.  A separate verifier compares the circuit
with multiplication in the eight-dimensional algebra generated by
$r,f,\tau$ on $48$ signed assignments.  Two modular line restrictions
give degrees $1600$ and $1609$, agreeing with the uniform proof above.
These are finite arithmetic audits; neither they nor the machine's formal
endpoint constitute a proof-assistant formalization of this new reduction.

The main remaining opportunities are to lower the degree of the
coprimality-enforcing multiplier restriction, to replace the three norms
by a less expensive elimination, and to combine the twelve-variable
construction with the sharper signed-domain analysis.  Formalizing the
branch separation and denominator argument would isolate a small reusable
lemma from the longer inherited native compiler proof.
